\chapter{Introduction}
\label{chap:intro}

\section{Survival Prediction in Healthcare}
\label{background}

Clinical prognosis concerns both whether an event occurs and when it occurs. A patient may remain free of recurrence for several years, experience an event soon after diagnosis, or reach the end of observation before the event is recorded. Survival analysis uses event timing and censored follow-up to represent these different observations. It provides a common framework for studying mortality, recurrence and other clinical time-to-event outcomes \cite{ravi14,ravi15}.

Right censoring occurs when observation ends before the event is recorded. The patient's exact event time is unknown, but their observed event-free duration remains informative. Survival models retain that information through their likelihoods, and survival-specific evaluation measures account for incomplete outcome status at a prediction horizon. This treatment of follow-up makes survival analysis appropriate for healthcare cohorts with varying observation periods \cite{ravi3,ravi14}.

The range of available methods has expanded from descriptive survival curves and regression to tree ensembles, neural networks and Bayesian models. Cox proportional-hazards regression provides a widely used model with directly interpretable coefficients. Random survival forest (RSF) represents nonlinear associations through ensembles of survival trees. DeepSurv uses a neural risk function within a proportional-hazards framework \cite{cox1972,ishwaran2008,ravi4}. Each offers a different balance of predictive flexibility, estimation and interpretation.

Bayesian survival modelling adds a distributional account of uncertainty. Prior distributions describe plausible parameter values, and the likelihood updates them using observed events and follow-up. Posterior distributions then support coefficient estimates, hazard ratios, credible intervals and probabilities of direction. These quantities are useful when several predictors contain overlapping information or the number of observed events is modest relative to the model's complexity \cite{ravi22,ibrahim2001}.

\section{Research Motivation}
\label{problem-statement}

The coefficient prior is a substantive part of a Bayesian survival model. A Normal prior applies a common scale to regression effects. A hierarchical shrinkage prior allows the amount of regularisation to vary across coefficients, while a continuous spike-and-slab specification interpolates between narrow and broad coefficient scales. These choices can alter posterior effect estimates, uncertainty and predictions even when the observed data and likelihood are identical \cite{carvalho2010,piironen2017,ravi25}.

The amount of survival information also varies. Sample size, observed event count, censoring and predictor dependence jointly influence how well effects can be estimated. A prior that gives useful regularisation in one information setting may behave differently after reducing the training sample or shortening follow-up. Understanding this interaction requires a comparison in which the likelihood, predictors and validation subjects remain common across prior specifications.

Interpretation presents a related question. A coefficient distribution describes an association and its uncertainty. A feature-contribution ranking describes which predictors contribute most strongly to a fitted risk function. SHAP attributes that function to predictors relative to a reference background \cite{lundberg2017,ravi30}. These summaries answer complementary questions. Their agreement and repeatability provide additional information about the explanation of a model beyond its prediction error.

\section{Research Contribution}
\label{research-gap}

This thesis contributes a survival benchmark centred on coefficient priors. Normal, horseshoe and continuous spike-and-slab priors share a full piecewise-exponential likelihood, baseline-hazard structure, predictors and validation partitions. Their comparison therefore concerns coefficient-prior choice within a defined survival model. CoxPH, RSF and DeepSurv provide reference predictions, while matched reductions in training information connect the prior differences to the evidence available for estimation.

The second contribution connects this predictive comparison to posterior interpretation. Coefficients and hazard ratios describe uncertain associations; grouped contributions and exact linear SHAP describe participation in the fitted risk function. Comparing their agreement and sampling stability distinguishes these explanatory properties empirically. Table~\ref{tab:closest-work} positions the combined design alongside the closest benchmarking, prior-modelling and explanation studies.

\section{Aim and Objectives}
\label{aim}

The aim is to assess how Normal, horseshoe and continuous spike-and-slab coefficient priors affect prediction, posterior inference and explanation stability in Bayesian piecewise-exponential proportional-hazards models for healthcare survival data.

The objectives are to:

\begin{enumerate}
\item construct a paired clinical survival benchmark with shared data preparation and validation;
\item quantify predictive accuracy, calibration and computational acceptance under the three coefficient priors;
\item interpret clinical associations and compare posterior contribution rankings with SHAP;
\item characterise how prior contrasts change under controlled information and modelling conditions.
\end{enumerate}

\section{Research Questions}
\label{research-questions}

\subsection*{RQ1: Prediction under different coefficient priors}
How do Normal, horseshoe and continuous spike-and-slab priors affect out-of-sample survival prediction, and how do their predictions compare with CoxPH, RSF and DeepSurv?

The primary predictive quantity is the within-cohort paired difference in inverse-probability-of-censoring-weighted integrated Brier score (IPCW IBS). Three pairwise prior contrasts in each cohort form a nine-comparison primary family. Discrimination, calibration and comparisons with the other model families provide complementary estimates.

\subsection*{RQ2: Information-regime robustness}
How does the prior effect change as training information falls or administrative censoring increases?

Matched training subsets and training-derived follow-up truncation define the information conditions. Differences-in-differences measure how a prior contrast changes between each reduced-information condition and its full-information reference.

\subsection*{RQ3: Explanation agreement and stability}
How stable are Bayesian global feature rankings, and how do they agree with SHAP rankings of the same fitted risk function?

Posterior expected absolute grouped contributions provide a common intrinsic importance measure across priors. They are compared with SHAP of the Bayesian posterior-mean log-risk function and, secondarily, with CoxPH explanations on the same subjects and reference backgrounds. Correlation, top-five overlap, rank variability and resampling quantify agreement and stability.

\subsection*{RQ4: Design robustness}
How do the main contrasts respond to choices about predictors, missingness, censoring, baseline-hazard intervals and prior settings?

The analysis compares effect direction, uncertainty and computational acceptance under the defined sensitivity conditions. Together, the four questions connect prior specification to prediction, interpretation and the conditions under which the fitted models are evaluated.

\section{Study Scope}
\label{scope-and-boundaries}

The study analyses GBSG2, METABRIC and WHAS500 separately, retaining each cohort's endpoint, time unit and clinical predictor definition. Repeated outer cross-validation evaluates predictions, with preprocessing and model selection fitted within training partitions. The Bayesian arms share six baseline-hazard intervals and differ in their coefficient priors.

Evaluation combines IPCW IBS, Uno's and Harrell's concordance, cumulative/dynamic AUC and calibration. Interpretation combines posterior associations with grouped contribution rankings and exact interventional linear SHAP of the posterior-mean log-risk. Original comparisons and prespecified sensitivities form the main analysis; the higher-target-acceptance horseshoe sensitivity and Holm bound are identified separately as post-hoc analyses. The conclusions concern the evaluated prediction procedures and their behaviour within these clinical cohorts.

\section{Thesis Structure}
\label{thesis-structure}

Chapter~\ref{chap:lit} reviews survival models, coefficient priors, evaluation and explanation, and develops the rationale for a controlled comparison. Chapter~\ref{chap:methods} specifies the cohorts, modelling framework, validation, diagnostics and sensitivity analyses. Chapter~\ref{chap:results} presents and discusses the predictive, posterior and robustness findings. Chapter~\ref{chap:conclusion} brings together the answers to the research questions and the study's methodological contribution.
